\section{Introduction} \label{sec:intro}

The Domain Name System (\ac{dns}) is a fundamental Internet component that
translates human-readable domain names into IP
addresses~\cite{RFC1034,RFC1035}. Whenever a user attempts to access a website
or online resource, a \ac{dns} query is issued to determine the corresponding
IP address. Despite its essential role, traditional \ac{dns} operates in
plaintext, making it vulnerable to eavesdropping, manipulation, and tracking.

Privacy-preserving protocols such as \ac{doh}~\cite{RFC8484} and
\ac{dot}~\cite{RFC7858} have been developed to address these concerns. These
protocols encrypt \ac{dns} traffic between the client and the resolver,
preventing on-path adversaries from observing or tampering with queries and
responses. While \ac{doh} and \ac{dot} significantly improve confidentiality,
they do not fully address privacy concerns -- especially when the resolver
becomes a point of surveillance or data aggregation.

To further decouple the identity of the client from the \ac{dns} query, more
advanced protocols like \ac{odns}~\cite{ODNS} and \ac{odoh}~\cite{RFC9230} have
been proposed. These protocols introduce cryptographic indirection layers,
ensuring no single server can correlate a query with its origin. However, these
schemes often rely on a limited number of specialized servers and assume that
parties do not collude -- a strong assumption in practice. Moreover, deployment
remains limited, and performance overhead can be substantial.

An alternative approach to achieving strong anonymity and unlinkability is to
route \ac{dns} queries through an anonymity network such as
Tor~\cite{dingledine2004}. Tor is a low-latency, onion-routing network designed
to anonymize Internet traffic by relaying it through multiple volunteer-run
nodes. In particular, Alec Muffett's proposal for \ac{dohot}~\cite{DoHoT}
leverages this infrastructure to anonymize \ac{dns} queries while maintaining
the advantages of encrypted transport.

Despite its privacy benefits, \ac{dohot} and similar approaches have
performance limitations. Tor circuits introduce additional latency due to
multiple hops all over the globe, and variability in relay performance can lead
to inconsistent results. As a result, users must often choose between strong
anonymity and usable performance.

In this study, we aim to explore whether it is possible to configure Tor to
mitigate these performance penalties without sacrificing the security and
privacy guarantees. To do so, we begin by defining the relevant threat models
and outlining the capabilities of potential passive and active adversaries in
the context of encrypted \ac{dns} transport.

Beyond the default Tor configuration, we evaluate the impact of three custom
circuit strategies aimed at improving latency:

\begin{enumerate}
    \item Constraining the EntryNode and/or ExitNode in the \texttt{torrc} file to
        be geographically close to the client.
    \item Using \texttt{carml}~\cite{carml} and \texttt{stem}~\cite{stem} to
        create three-hop and two-hop circuits with randomly selected relays.
    \item Creating three-hop and two-hop circuits with relays that are geographically
        proximate to the client.
\end{enumerate}

To eliminate the effects of caching and to maintain consistency across tests,
we issue queries for uniquely generated \acp{sld} under valid \acp{tld},
expected to return an NXDOMAIN \ac{rcode}. This approach ensures that the
observed latency mainly reflects the path between the client and the resolver
rather than the influence of cache hits or additional authoritative name servers.

Through these measurements, we aim to provide a fair and comprehensive
evaluation of various approaches to anonymizing \ac{dns} queries, with
particular attention to the trade-offs between privacy and performance.
Ultimately, we aim to identify configurations that offer practical improvements
for users who seek strong privacy guarantees without intolerable slowdowns.
